% !TeX root = ../../Thesis.tex
\ThesisChapter{\ThesisText{Implementation}{Implementierung}}{ch:implementation}
  {TODO: Add a brief chapter overview. \\
  This chapter describes the realization of the concept, the technologies used, and key implementation decisions.}

% TODO: Describe the implementation of the concepts and architecture presented in the previous chapter.

This chapter explains how the proposed solution was implemented in practice.
It should describe the relevant technical details required to understand and reproduce the implementation.

Possible contents include:
\begin{itemize}
  \item the selected programming languages, frameworks, libraries, and tools,
  \item the implementation of the individual components,
  \item the implementation of interfaces and communication mechanisms,
  \item relevant data structures, algorithms, and processing steps,
  \item configuration, deployment, and execution details, and
  \item important implementation decisions, challenges, and limitations.
\end{itemize}

Focus on details that are relevant to the thesis and necessary to understand the implementation.
Avoid describing trivial code or providing a complete source-code walkthrough.
Relevant code excerpts may be included where they help explain important implementation details.

The concepts and architecture chapter should describe the design of the solution, its components, responsibilities, and interactions.
This chapter should focus on the concrete technical realization of that design.

% Example for pseudocode typeset with algorithm2e.
\section{\ThesisText{Example Algorithm}{Beispiel: Algorithmus}}
\label{sec:example-algorithm}

Pseudocode is the right level of detail whenever the idea matters more than the syntax of a particular programming language.
Algorithm~\ref{alg:example-parallel-sum} states a block-wise parallel summation in this form.

\begin{algorithm}[htbp]
  \KwIn{Array $A$ with $n$ elements, thread count $p$}
  \KwOut{Sum of all elements of $A$}
  $\mathit{partial} \gets$ array of $p$ zeros\;
  \ForEach(\tcp*[f]{in parallel}){thread $t \in \{0, \dots, p-1\}$}{
    \For{$i \gets t \cdot n / p$ \KwTo $(t+1) \cdot n / p - 1$}{
      $\mathit{partial}[t] \gets \mathit{partial}[t] + A[i]$\;
    }
  }
  \Return $\sum_{t=0}^{p-1} \mathit{partial}[t]$\;
  \caption{\ThesisText{Block-wise parallel summation of an array.}{Blockweise parallele Summation eines Arrays.}}
  \label{alg:example-parallel-sum}
\end{algorithm}
